% SPDX-License-Identifier: MIT
% Copyright (c) 2026 Evgeny Baulin
% Optimization Methods, Seminar 05: one-page cheat sheet (English).
% Generated by checks/05. Second-Order and Quasi-Newton Methods.ipynb from web/05/js/content-cheatsheet.js and content/students.json. Do not edit by hand.
% Compile with: latexmk -xelatex -interaction=nonstopmode -halt-on-error Cheatsheet_en.tex
\RequirePackage{fix-cm}
\documentclass[10pt,a4paper]{article}

\usepackage{amsmath,amssymb}
\usepackage{fontspec}
\setmainfont{cmunrm.otf}[BoldFont=cmunbx.otf, ItalicFont=cmunti.otf, BoldItalicFont=cmunbi.otf]
\setsansfont{cmunss.otf}[BoldFont=cmunsx.otf, ItalicFont=cmunsi.otf, BoldItalicFont=cmunso.otf]
\setmonofont{cmuntt.otf}[BoldFont=cmuntb.otf, ItalicFont=cmunit.otf, BoldItalicFont=cmuntx.otf]
\usepackage{polyglossia}
\setdefaultlanguage{english}

\usepackage[a4paper,left=8mm,right=8mm,top=7mm,bottom=11mm,footskip=5.5mm]{geometry}
\usepackage{xcolor}
\usepackage{multicol}
\usepackage{enumitem}
\usepackage{ragged2e}
\usepackage{fancyhdr}
\usepackage[colorlinks=true,linkcolor=hseblue,urlcolor=hseblue]{hyperref}

\definecolor{hseblue}{RGB}{16,45,105}
\definecolor{hseaccent}{RGB}{198,40,40}
\definecolor{hsegray}{RGB}{108,114,126}
\definecolor{hselight}{RGB}{240,243,248}

\hypersetup{pdftitle={Seminar 05 cheat sheet: Second-Order and Quasi-Newton Methods}, pdfauthor={Evgeny Baulin}}

% Page style: the author and the license in the footer.
\pagestyle{fancy}
\fancyhf{}
\renewcommand{\headrulewidth}{0pt}
\fancyfoot[L]{\scriptsize\color{hsegray}\textcopyright\ 2026 Evgeny Baulin, \href{https://t.me/tarakan_tuc}{t.me/tarakan\_tuc}}
\fancyfoot[R]{\scriptsize\color{hsegray}MIT License}

% Compact layout.
\setlength{\parindent}{0pt}
\setlength{\parskip}{0pt}
\setlength{\columnsep}{4.5mm}
\setlength{\columnseprule}{0.3pt}
\renewcommand{\columnseprulecolor}{\color{hselight}}
\setlength{\emergencystretch}{2em}
\setlength{\RaggedRightRightskip}{0pt plus 4em}
\binoppenalty=2000
\relpenalty=500
\setlength{\abovedisplayskip}{2pt}
\setlength{\belowdisplayskip}{2pt}
\setlist[itemize]{leftmargin=7pt, labelsep=3pt, itemsep=1.2pt, topsep=1pt, parsep=0pt, partopsep=0pt,
  label={\color{hseaccent}\scriptsize\textbullet}}

\makeatletter
\newcommand{\cssec}[1]{\par\addvspace{3.5pt}\noindent
  \colorbox{hselight}{\parbox{\dimexpr\linewidth-2\fboxsep}{\RaggedRight\hyphenpenalty=10000\color{hseblue}\bfseries\strut #1}}%
  \par\nobreak\vspace{1pt}\@afterheading}
\makeatother
\clubpenalty=10000
\widowpenalty=10000

\begin{document}
\thispagestyle{fancy}

\noindent
\begin{minipage}[b]{0.6\textwidth}
  \RaggedRight\hyphenpenalty=10000
  {\color{hseaccent}\bfseries\small Seminar 05 cheat sheet}\\[1pt]
  {\color{hseblue}\bfseries\Large Second-Order\\ and Quasi-Newton Methods}
\end{minipage}\hfill
\begin{minipage}[b]{0.39\textwidth}
  \raggedleft\footnotesize
  {\bfseries Optimization Methods}\\
  HSE University, Faculty of Computer Science, 2026--2027\\
  Course lecturer: Andrey Ignatov\\
  Author: Evgeny Baulin, Telegram \href{https://t.me/tarakan_tuc}{@tarakan\_tuc}
\end{minipage}

\vspace{2pt}
{\color{hseblue}\rule{\textwidth}{0.8pt}}
\vspace{-4pt}

\fontsize{9}{11.4}\selectfont
\begin{multicols}{3}
\RaggedRight


% ---- newton
\cssec{Newton step}
\begin{itemize}
  \item $m(p)=f+g^\top p+\tfrac{1}{2}p^\top Hp$; $Hp_N=-g$, $x^+=x+p_N$.
  \item Solve by Cholesky or LU, never with $H^{-1}$.
  \item Quadratic, $Q\succ0$: one step, $x^+=Q^{-1}b$.
  \item $\nabla^2 f(x^*)\succ0$, $M$-Lipschitz $H$: $\|e^+\|\le\tfrac{M}{2m}\|e\|^2$.
  \item Singular $H(x^*)$: only linear, e.g. $x^4$ gives $x^+=\tfrac{2}{3}x$.
\end{itemize}

% ---- fail
\cssec{Failure modes}
\begin{itemize}
  \item $g^\top p_N=-g^\top H^{-1}g$: may be $>0$ if $H$ is indefinite.
  \item Newton solves $\nabla f=0$: saddles and maxima attract.
  \item Good model at $x$, step outside its region: $f$ may jump up.
\end{itemize}

% ---- guard
\cssec{Safeguards}
\begin{itemize}
  \item Armijo: $f(x+\alpha p)\le f+c\alpha g^\top p$, $\alpha=1,\beta,\beta^2,\dots$
  \item $c<\tfrac{1}{2}$, $H(x^*)\succ0$: eventually $\alpha=1$, quadratic rate kept.
  \item $(H+\lambda I)p=-g$, $\lambda>-\lambda_{\min}(H)$; raise $\lambda$ until Cholesky works.
  \item $\lambda\to0$: Newton; $\lambda\to\infty$: $p\approx-g/\lambda$.
\end{itemize}

% ---- ls
\cssec{Least squares}
\begin{itemize}
  \item $f=\tfrac{1}{2}\|r\|^2$: $\nabla f=J^\top r$, $\nabla^2 f=J^\top J+\sum r_i\nabla^2 r_i$.
  \item Gauss–Newton $J^\top Jp=-J^\top r$; LM $(J^\top J+\lambda I)p=-J^\top r$.
\end{itemize}

% ---- cost
\cssec{Cost}
\begin{itemize}
  \item Dense: $O(n^2)$ memory, $\approx n^3/3$ for Cholesky.
  \item Logistic: $Hv=N^{-1}X^\top(d\odot Xv)+\gamma v$, $d=\sigma(1-\sigma)$.
  \item Newton-CG: $k$ products cost $\approx2kNn$ vs $Nn^2$ to form $H$.
\end{itemize}

% ---- bfgs
\cssec{BFGS}
\begin{itemize}
  \item $s=x^+-x$, $y=g^+-g$; $B^+s=y$, $G^+y=s$.
  \item $G^+=(I-\rho sy^\top)G(I-\rho ys^\top)+\rho ss^\top$, $\rho=1/y^\top s$.
  \item $B^+=B-\tfrac{Bss^\top B}{s^\top Bs}+\tfrac{yy^\top}{y^\top s}$; $p=-Gg$.
  \item $G\succ0$, $y^\top s>0$ $\Rightarrow$ $G^+\succ0$. Superlinear near $x^*$.
  \item Exact line search on an $n$-dim quadratic: at most $n$ steps.
\end{itemize}

% ---- curv
\cssec{Curvature condition}
\begin{itemize}
  \item $y^\top s$ = average curvature along $s$; $\ge\mu\|s\|^2$ if strongly convex.
  \item Wolfe: $\nabla f(x+\alpha p)^\top p\ge c_2g^\top p$ $\Rightarrow$ $y^\top s>0$.
  \item Skip if $y^\top s\le\varepsilon\|y\|\|s\|$, or damp.
\end{itemize}

% ---- lbfgs
\cssec{L-BFGS}
\begin{itemize}
  \item Keep $m$ pairs; $q=g$; newest to oldest: $a_i=\rho_is_i^\top q$, $q\mathrel{-}=a_iy_i$.
  \item $r=\gamma q$, $\gamma=\tfrac{s^\top y}{y^\top y}$ (last pair).
  \item Oldest to newest: $b_i=\rho_iy_i^\top r$, $r\mathrel{+}=s_i(a_i-b_i)$; $p=-r$.
  \item $\approx4mn$ flops, $2mn$ numbers. Bounds: L-BFGS-B, not clipping.
\end{itemize}

\end{multicols}
\end{document}
